\titledquestion{Machine Learning \& Neural Networks}[8] 
\begin{parts}

    
    \part[4] Adam Optimizer\newline
        Recall the standard Stochastic Gradient Descent update rule:
        \alns{
            	\btheta_{t+1} &\gets \btheta_t - \alpha \nabla_{\btheta_t} J_{\text{minibatch}}(\btheta_t)
        }
        where $t+1$ is the current timestep, $\btheta$ is a vector containing all of the model parameters, ($\btheta_t$ is the model parameter at time step $t$, and $\btheta_{t+1}$ is the model parameter at time step $t+1$), $J$ is the loss function, $\nabla_{\btheta} J_{\text{minibatch}}(\btheta)$ is the gradient of the loss function with respect to the parameters on a minibatch of data, and $\alpha$ is the learning rate.
        Adam Optimization\footnote{Kingma and Ba, 2015, \url{https://arxiv.org/pdf/1412.6980.pdf}} uses a more sophisticated update rule with two additional steps.\footnote{The actual Adam update uses a few additional tricks that are less important, but we won't worry about them here. If you want to learn more about it, you can take a look at: \url{http://cs231n.github.io/neural-networks-3/\#sgd}}
            
        \begin{subparts}

            \subpart[2]First, Adam uses a trick called {\it momentum} by keeping track of $\bm$, a rolling average of the gradients:
                \alns{
                	\bm_{t+1} &\gets \beta_1\bm_{t} + (1 - \beta_1)\nabla_{\btheta_t} J_{\text{minibatch}}(\btheta_t) \\
                	\btheta_{t+1} &\gets \btheta_t - \alpha \bm_{t+1}
                }
                where $\beta_1$ is a hyperparameter between 0 and 1 (often set to  0.9). Briefly explain in 2--4 sentences (you don't need to prove mathematically, just give an intuition) how using $\bm$ stops the updates from varying as much and why this low variance may be helpful to learning, overall.\newline
              
                
            \subpart[2] Adam extends the idea of {\it momentum} with the trick of {\it adaptive learning rates} by keeping track of  $\bv$, a rolling average of the magnitudes of the gradients:
                \alns{
                	\bm_{t+1} &\gets \beta_1\bm_{t} + (1 - \beta_1)\nabla_{\btheta_t} J_{\text{minibatch}}(\btheta_t) \\
                	\bv_{t+1} &\gets \beta_2\bv_{t} + (1 - \beta_2) (\nabla_{\btheta_t} J_{\text{minibatch}}(\btheta_t) \odot \nabla_{\btheta_t} J_{\text{minibatch}}(\btheta_t)) \\
                	\btheta_{t+1} &\gets \btheta_t - \alpha \bm_{t+1} / \sqrt{\bv_{t+1}}
                }
                where $\odot$ and $/$ denote elementwise multiplication and division (so $\bz \odot \bz$ is elementwise squaring) and $\beta_2$ is a hyperparameter between 0 and 1 (often set to  0.99). Since Adam divides the update by $\sqrt{\bv}$, which of the model parameters will get larger updates?  Why might this help with learning?
                
           
                
                \end{subparts}
                \ifans{
                    \begin{enumerate}[label=\roman*.]
                        \item 通过累积过去的梯度值，会让当前的参数变化更加平滑，这种方法能够减少因噪声梯度而导致的参数变化，
                        并减少梯度方向波动时的震荡程度，momentum有助于保持参数更新的速度，防止学习陷入平台期或局部最小值。
                        \item 在梯度中存在较小的非中心值的参数会得到更大的更新。除了使用一成不变的学习率，$/ \sqrt{\bv}$使得学习率能够适应不同的参数，
                        这会让参数更新更加平滑，避免一下子更新过多错过了最优点，或者让每次梯度变化小的参数能够更快地更新。
                    \end{enumerate}
                }
        
        
            \part[4] 
            Dropout\footnote{Srivastava et al., 2014, \url{https://www.cs.toronto.edu/~hinton/absps/JMLRdropout.pdf}} is a regularization technique. During training, dropout randomly sets units in the hidden layer $\bh$ to zero with probability $p_{\text{drop}}$ (dropping different units each minibatch), and then multiplies $\bh$ by a constant $\gamma$. We can write this as:
                \alns{
                	\bh_{\text{drop}} = \gamma \bd \odot \bh
                }
                where $\bd \in \{0, 1\}^{D_h}$ ($D_h$ is the size of $\bh$)
                is a mask vector where each entry is 0 with probability $p_{\text{drop}}$ and 1 with probability $(1 - p_{\text{drop}})$. $\gamma$ is chosen such that the expected value of $\bh_{\text{drop}}$ is $\bh$:
                \alns{
                	\mathbb{E}_{p_{\text{drop}}}[\bh_\text{drop}]_i = h_i \text{\phantom{aaaa}}
                }
                for all $i \in \{1,\dots,D_h\}$. 
            \begin{subparts}
            \subpart[2]
                What must $\gamma$ equal in terms of $p_{\text{drop}}$? Briefly justify your answer or show your math derivation using the equations given above.
                
            
          \subpart[2] Why should dropout be applied during training? Why should dropout \textbf{NOT} be applied during evaluation? \textbf{Hint:} it may help to look at the dropout paper linked. \newline
          
          \ifans{
                \begin{enumerate}[label=\roman*.]
                    \item 
                    \begin{align}
                        \mathbb{E}_{p_{\text{drop}}}[\bh_\text{drop}]_i 
                        &= \mathbb{E}_{p_{\text{drop}}}[\gamma \bd \odot \bh_{i}] \\
                        &= \gamma \bh_{i} \mathbb{E}_{p_\text{drop}}[\bd] \\
                        &= \gamma \bh_{i} [0 \times p_{\text{drop}} + 1 \times (1 - p_{\text{drop}})] \\
                        &= \gamma \bh_{i} (1 - p_{\text{drop}}) \\
                        &= \bh_{i}
                    \end{align}
                    $$
                        \gamma \bh_{i} (1 - p_{\text{drop}}) = \bh_{i}
                    $$
                    $$
                        \gamma = \frac{1}{1 - p_{\text{drop}}}
                    $$
                    要让$\bh_{\text{drop}}$的期望等于$\bh$，那么根据前面的描述我们知道，某神经元被drop的概率是$p_{\text{drop}}$，也就是说
                    在该位置为0的概率也就是$p_{\text{drop}}$，同理可得该位置为1的概率，那么就可以得到这个位置的期望$\mathbb{E}_{p_{\text{drop}}}[\cdot]_{i}$。
                    \item 在训练过程中使用dropout可以防止过拟合，避免在训练过程中过度依赖于某些神经元的学习，
                    通过dropout来消去一部分神经元的作用，可以让其他神经元更好地学习数据的特征，同时也训练了更多不同的子神经网络，
                    使得神经网络具有更强地学习能力和鲁棒性。
                    
                    在评估的时候不使用dropout，是为了在评估的时候能够充分利用神经网络的能力，从其生成的结果和标准答案之间的不同，来判断当前神经网络的参数效果如何
                    测试的时候不需要进行额外的操作，可以加快推理速度。
                \end{enumerate}
          }
         
        \end{subparts}


\end{parts}
